\section{代数曲线基础：只讲够用的语言}
\label{sec:ch07-algebraic-varieties}
% ============================================================================

一旦我们决定把模曲线看成定义在 $\Q$ 上的代数曲线，
立刻会遇到一个很现实的问题：
前面几章的语言几乎全是复分析与黎曼面，
而“定义在 $\Q$ 上”“模 $p$ 约化”“Galois 作用”这些句子却属于代数几何。
如果没有新的词汇，我们根本无法准确地说出
“同一条复曲线其实来自 $\Q$ 上的方程”到底是什么意思。

\textbf{我们遇到了什么问题？}
我们已经会在黎曼面上谈全纯函数、亚纯函数、除子与 Riemann--Roch；
但现在需要把这些对象改写成不依赖复解析坐标的代数语言。
特别是，当底域从 $\C$ 变成 $\Q$ 或 $\F_p$ 时，
“全纯”这个词已经没有意义，
可“函数域”“有理映射”“除子”仍然必须存在。
所以本节的任务不是系统学习代数几何，
而是提炼出之后讨论模曲线所必需的最小工具箱。

\textbf{核心思想是什么？}
核心思想是：
对一条光滑射影曲线来说，
复分析里最重要的那些对象都可以完全代数化。
紧黎曼面上的亚纯函数变成函数域 $K(C)$ 中的元素；
零点极点的计数变成除子 $\Div(C)$；
“允许在给定点有若干极”的条件变成向量空间 $\mathcal L(D)$；
而第\ref{ch:dimension}章的 Riemann--Roch 定理则可以原封不动地翻译到这个框架中。

\textbf{引入之后能得到什么？}
有了这套语言，
我们就能在 $\Q$ 上谈曲线、在 $\F_p$ 上谈约化、
在不同模型之间谈有理映射，
并且仍保留前几章已经练熟的除子与维数工具。
特别地，
当我们说“$X_0(N)$ 有一个 $\Q$-模型”时，
意思将不再是含糊的口号，
而是“存在一条光滑射影曲线 $C/\Q$，其函数域、除子与映射都在 $\Q$ 上定义”。

\begin{figure}[ht]
\centering
\begin{tikzpicture}[>=Stealth, font=\small]
  \draw[thick] (-3.2,-2.1) -- (-1.1,-2.9) -- (-0.1,-0.8) -- cycle;
  \node at (-2.2,-3.2) {$\mathbb P^2$};
  \draw[thick, blue!70!black] plot[smooth cycle] coordinates {(-2.7,-1.9) (-2.0,-1.0) (-1.2,-1.4) (-1.0,-2.2) (-2.0,-2.5)};
  \node[blue!70!black] at (-2.1,-0.5) {$C$};

  \draw[thick] (1.1,-2.1) -- (3.2,-2.9) -- (4.2,-0.8) -- cycle;
  \node at (2.1,-3.2) {$\mathbb P^2$};
  \draw[thick, green!50!black] plot[smooth cycle] coordinates {(1.5,-1.8) (2.3,-1.0) (3.2,-1.2) (3.4,-2.2) (2.4,-2.6)};
  \node[green!50!black] at (2.6,-0.5) {$D$};

  \draw[->, very thick] (-0.1,-1.8) -- node[above] {$\varphi$} (1.1,-1.8);
  \node at (0.5,-2.4) {有限有理映射};
\end{tikzpicture}
\caption{本节只保留最小必需的代数几何语言：把曲线看成射影空间中的方程，再研究曲线之间的有理映射、函数域与除子。}
\label{fig:ch07-algebraic-varieties}
\end{figure}


\begin{definition}[射影空间与射影曲线]
\label{def:ch07-projective-curve}
设 $K$ 是一域。
$n$ 维\textbf{射影空间}\index{射影空间!projective space} $\mathbb P^n_K$ 定义为
$K^{n+1}\setminus\{0\}$ 按比例关系
\[
  (x_0,\dots,x_n)\sim (\lambda x_0,\dots,\lambda x_n)\qquad (\lambda\in K^\times)
\]
取商所得的集合，其点记为齐次坐标 $[x_0:\cdots:x_n]$。
若若干齐次多项式 $F_1,\dots,F_r\in K[X_0,\dots,X_n]$ 的公共零点集
\[
  C=V(F_1,\dots,F_r)\subset \mathbb P^n_K
\]
是一维代数簇，
则称 $C$ 为一条\textbf{射影代数曲线}\index{代数曲线!射影代数曲线}。
若在每个点上 Jacobian 矩阵满秩，则称 $C$ 光滑。
\end{definition}

\begin{example}[椭圆曲线作为平面射影曲线]
\label{ex:ch07-elliptic-projective}
设 $a,b\in \Q$，并满足
\[
  \Delta=-16(4a^3+27b^2)\neq 0.
\]
则仿射方程
\[
  y^2=x^3+ax+b
\]
的射影闭包
\[
  Y^2Z=X^3+aXZ^2+bZ^3\subset \mathbb P^2_{\Q}
\]
是一条光滑曲线。
它在无穷远处只有一点 $[0:1:0]$，
这就是最熟悉的椭圆曲线模型。
\end{example}

\begin{proof}
在仿射片 $Z=1$ 上，
方程变为 $F(x,y)=y^2-x^3-ax-b=0$。
若存在奇点，则必有
\[
  F=0,\qquad \frac{\partial F}{\partial x}=-3x^2-a=0,
  \qquad \frac{\partial F}{\partial y}=2y=0.
\]
由 $2y=0$ 得 $y=0$，再由 $-3x^2-a=0$ 得 $x$ 是三次多项式 $x^3+ax+b$ 的重根。
这等价于其判别式 $-4a^3-27b^2$ 为零，
与 $\Delta\neq 0$ 矛盾。
所以仿射部分无奇点。

无穷远点只能满足 $Z=0$，此时方程化为 $X^3=0$，故 $X=0$，只剩 $[0:1:0]$。
在该点附近取局部坐标 $u=X/Y$、$v=Z/Y$，
方程变为
\[
  v=u^3+auv^2+bv^3.
\]
于是
\[
  \frac{\partial}{\partial v}\bigl(v-u^3-auv^2-bv^3\bigr)\Big|_{(0,0)}=1\neq 0,
\]
故此点也光滑。
因此整条射影曲线光滑。
\end{proof}

\begin{definition}[函数域与有理函数]
\label{def:ch07-function-field}
设 $C/K$ 是一条不可约光滑射影曲线。
它的\textbf{函数域}\index{函数域!curve} $K(C)$ 定义为 $C$ 上有理函数所成的域。
也就是说，$K(C)$ 的元素可以在某个非空 Zariski 开集上写成正则函数之比。
\end{definition}

对曲线来说，函数域是一条极其高效的编码。
许多几何性质都可以直接从 $K(C)$ 读出：
点给出赋值，映射给出域嵌入，
而除子则记录函数在各点的零极行为。

\begin{definition}[除子]
\label{def:ch07-divisor}
设 $C/K$ 为光滑射影曲线。
其\textbf{除子群}\index{除子!curve}定义为
\[
  \Div(C)=\bigoplus_{P\in C}\Z\cdot P.
\]
也就是说，一个除子是有限和
\[
  D=\sum_P n_P P,\qquad n_P\in\Z.
\]
其\textbf{次数}定义为
\[
  \deg D=\sum_P n_P.
\]
若 $f\in K(C)^\times$，记 $\ord_P(f)$ 为 $f$ 在 $P$ 处的零极阶，
则
\[
  \divisor(f)=\sum_P \ord_P(f)P
\]
称为 $f$ 的\textbf{主除子}\index{主除子!principal divisor}。
\end{definition}

\begin{proposition}[主除子的次数为零]
\label{prop:ch07-principal-degree-zero}
对任意 $f\in K(C)^\times$，都有
\[
  \deg \divisor(f)=0.
\]
\end{proposition}

\begin{proof}
把 $f$ 看成到射影直线的有理映射
\[
  f:C\longrightarrow \mathbb P^1,
  \qquad P\longmapsto [f(P):1],
\]
并在极点处理解为送到 $\infty=[1:0]$。
因为 $C$ 光滑射影，非零有理函数必给出到 $\mathbb P^1$ 的非恒定态射。
设其次数为 $d$。

$f$ 的零点恰是 $f^{-1}(0)$ 中的点，按重数计数得到
\[
  \sum_{\ord_P(f)>0} \ord_P(f)=d,
\]
因为 $0\in \mathbb P^1$ 的原像总重数等于映射次数。
同理，极点恰是 $f^{-1}(\infty)$ 中的点，按重数计数也得到
\[
  \sum_{\ord_P(f)<0} -\ord_P(f)=d.
\]
于是
\[
  \deg\divisor(f)
  =\sum_{\ord_P(f)>0}\ord_P(f)-\sum_{\ord_P(f)<0}(-\ord_P(f))
  =d-d=0.
\]
证毕。
\end{proof}

\begin{definition}[Riemann--Roch 空间]
\label{def:ch07-l-of-d}
对除子 $D=\sum n_PP$，定义
\[
  \mathcal L(D)=\{f\in K(C)^\times:\divisor(f)+D\ge 0\}\cup\{0\},
\]
即所有在每个点 $P$ 处至多有 $n_P$ 阶极的有理函数所成的 $K$-向量空间。
记
\[
  \ell(D)=\dim_K \mathcal L(D).
\]
\end{definition}

这个记号与第\ref{ch:dimension}章完全平行：
那里我们用亚纯函数和除子定义同样的空间，
这里则把一切写成纯代数语言。
因此第\ref{ch:dimension}章证明过的很多直觉，
现在都可以直接照搬过来。

\begin{theorem}[代数形式的 Riemann--Roch]
\label{thm:ch07-algebraic-riemann-roch}
设 $C/\C$ 是亏格为 $g$ 的光滑射影曲线，$K_C$ 为其典范除子。
则对任意除子 $D$，有
\[
  \ell(D)-\ell(K_C-D)=\deg D+1-g.
\]
\end{theorem}

\begin{proof}
由第\ref{ch:riemann-surfaces}章与 GAGA 视角，
$C(\C)$ 可以看成一条紧黎曼面 $X$。
在这条紧黎曼面上，除子、亚纯函数与全纯微分的概念与代数定义完全对应：
$\mathcal L(D)$ 就是第\ref{ch:dimension}章里记作 $L(D)$ 的亚纯函数空间，
而典范除子 $K_C$ 对应于任一非零亚纯微分的除子类。
因此本定理左边的两项，
正是第\ref{ch:dimension}章 \cref{thm:riemann-roch} 中出现的两个维数。

更具体地说，
$\ell(D)$ 是满足 $\divisor(f)+D\ge 0$ 的亚纯函数维数；
$\ell(K_C-D)$ 则等于满足
\[
  \divisor(\omega)\ge D
\]
的亚纯微分空间维数，
这正是解析版 Riemann--Roch 中的 $i(D)$。
于是直接把 \cref{thm:riemann-roch} 应用于紧黎曼面 $X$，便得到
\[
  \ell(D)-\ell(K_C-D)=\deg D+1-g.
\]
这正是所求公式。
\end{proof}

\textbf{注。}
本章并不需要更高维概型上的 Riemann--Roch，
也不需要层上同调的完整形式。
对模曲线而言，一维情形已经足够强大：
它既能控制模形式空间的维数，
也能控制之后 Hecke 对应与 Jacobian 上的代数操作。

\begin{definition}[有理映射与定义在 $\Q$ 上的态射]
\label{def:ch07-morphism-over-q}
设 $C,D$ 是定义在 $\Q$ 上的光滑射影曲线。
一个\textbf{有理映射}\index{有理映射!curve} $\varphi:C\dashrightarrow D$，
是指在 $C$ 的某个非空开集上定义的态射。
若它在整个 $C$ 上都有定义，
且局部表达式的系数都落在 $\Q$ 中，
则称 $\varphi$ 是一条\textbf{定义在 $\Q$ 上的态射}。
\end{definition}

\begin{proposition}[曲线上的非恒定态射与函数域]
\label{prop:ch07-morphism-function-field}
设 $\varphi:C\to D$ 是域 $K$ 上两条光滑射影曲线之间的非恒定态射。
则拉回给出单射
\[
  \varphi^*:K(D)\hookrightarrow K(C),
\]
且
\[
  \deg \varphi=[K(C):\varphi^*K(D)].
\]
特别地，$\varphi$ 是有限态射。
\end{proposition}

\begin{proof}
若 $f\in K(D)$ 且 $\varphi^*f=0$，
则 $f\circ\varphi$ 在 $C$ 上恒为零。
因为 $\varphi$ 非恒定，其像在不可约曲线 $D$ 中稠密，
所以 $f$ 只能是零函数。
故 $\varphi^*$ 单射。

接下来证明次数公式。
取 $D$ 上任一非空仿射开集 $U$，其坐标环记为 $A$；
再取 $V=\varphi^{-1}(U)$ 中的非空仿射开集，坐标环记为 $B$。
由于 $\varphi$ 非恒定且曲线是一维，
$B$ 是 $A$ 的整扩张，故诱导的函数域扩张 $K(C)/K(D)$ 有有限次数。
这说明 $\varphi$ 有有限纤维，也就是有限态射。

对 $Q\in D$ 的一般点，纤维
\[
  \varphi^{-1}(Q)=\sum_{P\mapsto Q} e_P\,P
\]
的总重数与域扩张次数相等。
这是曲线上一维赋值理论的基本事实：
每个点 $P$ 给出 $K(C)$ 中延拓 $K(D)$ 上点 $Q$ 的赋值，
其分歧指数为 $e_P$，
而一般点上没有剩余扩张，于是
\[
  [K(C):K(D)]=\sum_{P\mapsto Q} e_P.
\]
右边正是态射次数的定义。
故
\[
  \deg\varphi=[K(C):\varphi^*K(D)].
\]
证毕。
\end{proof}

\begin{example}[定义在 $\Q$ 上的椭圆曲线]
\label{ex:ch07-elliptic-over-q}
当 $a,b\in\Q$ 且 $4a^3+27b^2\neq 0$ 时，
\[
  E: y^2=x^3+ax+b
\]
定义了一条 $\Q$ 上的光滑射影曲线。
它的函数域是 $\Q(E)=\Q(x,y)/(y^2-x^3-ax-b)$，
零元由无穷远点给出。
因此“椭圆曲线 $E/\Q$”本质上就是“带一个 $\Q$-有理点的亏格 $1$ 曲线”。
这正是本章要赋予 $X_0(N)$ 的那种代数身份。
\end{example}

到这里为止，我们已经有了讨论模曲线所需的基本语法：
方程、函数域、除子、Riemann--Roch 与曲线之间的态射。
下一节就把这些工具全部用到 $X_0(N)$ 身上，
先用 $j$-函数写出显式代数方程，
再用模解释说明这条方程为什么天然定义在 $\Q$ 上。
